\documentclass[12pt,a4paper]{article}

% ---------- Packages ----------
\usepackage[margin=1in]{geometry}
\usepackage{graphicx}
\usepackage{float}
\usepackage{hyperref}
\usepackage{booktabs}
\usepackage{enumitem}
\usepackage{xcolor}
\usepackage{listings}
\usepackage{amsmath}
\usepackage{amssymb}

% ---------- Hyperref ----------
\hypersetup{
  colorlinks=true,
  linkcolor=blue,
  urlcolor=blue,
  citecolor=blue
}

% ---------- Code listing style ----------
\lstdefinestyle{cstyle}{
  language=C,
  basicstyle=\ttfamily\small,
  keywordstyle=\color{blue},
  commentstyle=\color{gray},
  stringstyle=\color{teal},
  numbers=left,
  numberstyle=\tiny\color{gray},
  stepnumber=1,
  numbersep=6pt,
  showstringspaces=false,
  breaklines=true,
  frame=single,
  tabsize=2
}

\title{\textbf{Webots Autonomous Exploration and Target Navigation Report}\\
\large (A* Path Planning + Frontier Exploration + Color-Pillar Targeting + Safety Manager)}

\begin{document}
\maketitle

\tableofcontents
\newpage

% =========================================================
% =========================================================
\section{System Overview (Main Control Loop)}
\textbf{Key functions:} \texttt{main()}, \texttt{while(wb\_robot\_step(...))}, \texttt{switch(state)}

\begin{enumerate}[leftmargin=2em]
  \item \textbf{Safety Update:} reads IR sensors and pitch, possibly stops motion and enters recovery.
  \item \textbf{Mapping Update:} integrates LiDAR scans into an occupancy grid.
  \item \textbf{Perception Update:} marks additional obstacles from camera-based green/red detection.
  \item \textbf{Cost Map Update:} generates inflated costs based on obstacle distance (adaptive inflation).
  \item \textbf{Visualization:} renders the map, frontiers, path, and safety status.
  \item \textbf{State Machine:} selects between frontier exploration and target navigation (blue then yellow).
\end{enumerate}

% =========================================================
\section{Data Structures and Global Maps}
\textbf{Key functions:} \texttt{init\_grid()}, \texttt{world\_to\_grid()}, \texttt{is\_valid\_cell()}

\subsection{Grid Representation}
The environment is discretized into a 2D grid with resolution \texttt{GRID\_RESOLUTION} (meters per cell).
\begin{itemize}[leftmargin=2em]
  \item \texttt{grid[y][x]}: cell type (unknown/free/obstacle/robot/blue pillar/yellow pillar)
  \item \texttt{obstacle\_counter[y][x]} and \texttt{free\_counter[y][x]}: temporal filters for stability
  \item \texttt{cost\_map[y][x]}: numeric planning cost used by A*
  \item \texttt{frontier\_edge\_map[y][x]}: computed frontier edges for exploration
\end{itemize}

\subsection{Cell Types}
\begin{center}
\begin{tabular}{ll}
\toprule
\textbf{Constant} & \textbf{Meaning} \\
\midrule
\texttt{CELL\_UNKNOWN} & Unknown/unobserved area \\
\texttt{CELL\_FREE} & Traversable space \\
\texttt{CELL\_OBSTACLE} & Blocked cell (LiDAR/camera/safety marked) \\
\texttt{CELL\_BLUE\_PILLAR} & Blue target marker in grid \\
\texttt{CELL\_YELLOW\_PILLAR} & Yellow target marker in grid \\
\bottomrule
\end{tabular}
\end{center}

% =========================================================
\section{A* Path Planning Implementation}
\textbf{Key functions:} \texttt{ASPathCreate()}, \texttt{ASNeighborListAdd()}, \texttt{ASPathGetNode()},
\texttt{find\_path\_astar()}, \texttt{destroy\_grid\_path()}

\subsection{Generic A* Engine}
The controller contains a standalone A* implementation that is decoupled from Webots. The A* engine uses callbacks to:
\begin{itemize}[leftmargin=2em]
  \item fetch neighbor nodes,
  \item compute heuristic cost,
  \item compare node keys,
  \item optionally early-exit if too many nodes are expanded.
\end{itemize}

\subsection{Grid-Specific A* Callbacks}
\textbf{Key functions:} \texttt{grid\_node\_neighbors()}, \texttt{grid\_path\_cost\_heuristic()},
\texttt{grid\_node\_comparator()}, \texttt{grid\_early\_exit()}

\begin{itemize}[leftmargin=2em]
  \item \texttt{grid\_node\_neighbors()}: expands 8-connected neighbors. Each edge cost is scaled by the target cell's cost from \texttt{cost\_map}.
  \item \texttt{grid\_path\_cost\_heuristic()}: Euclidean distance (admissible heuristic in grid).
  \item \texttt{grid\_node\_comparator()}: deterministic ordering by \texttt{(y,x)}.
  \item \texttt{grid\_early\_exit()}: stops if visited nodes exceed a threshold (prevents long stalls).
\end{itemize}

\subsection{Planning}
\texttt{find\_path\_astar()} converts start/goal grid cells into an \texttt{ASPath} and then stores the resulting sequence into \texttt{GridPath} (array of waypoints).

% =========================================================
\section{Mapping with LiDAR}
\textbf{Key functions:} \texttt{process\_lidar()}, \texttt{draw\_line\_on\_grid()},
\texttt{decay\_counters()}, \texttt{filter\_connected\_components()},
\texttt{should\_integrate\_lidar()}

\subsection{LiDAR Integration}
\texttt{process\_lidar()} reads the LiDAR range image and for each ray:
\begin{enumerate}[leftmargin=2em]
  \item computes the ray angle relative to robot heading,
  \item projects the measured endpoint to world coordinates,
  \item converts world coordinates to grid coordinates,
  \item draws the ray using \texttt{draw\_line\_on\_grid()}.
\end{enumerate}

\subsection{Ray Tracing and Temporal Filtering}
\texttt{draw\_line\_on\_grid()} performs Bresenham-style ray tracing:
\begin{itemize}[leftmargin=2em]
  \item intermediate cells are reinforced as free (increase \texttt{free\_counter}),
  \item endpoint cell is reinforced as obstacle (increase \texttt{obstacle\_counter}),
  \item \textbf{pillar cells are protected} (LiDAR cannot overwrite \texttt{CELL\_BLUE\_PILLAR} or \texttt{CELL\_YELLOW\_PILLAR}).
\end{itemize}

\subsection{LiDAR Updates During Unstable Motion}
\texttt{should\_integrate\_lidar()} can disable LiDAR integration when yaw rate is high or pitch is abnormal. This reduces mapping corruption while rotating quickly or climbing.

\subsection{Noise Decay and Blob Filtering}
\begin{itemize}[leftmargin=2em]
  \item \texttt{decay\_counters()}: decreases confidence over time to remove stale false obstacles.
  \item \texttt{filter\_connected\_components()}: removes small isolated obstacle blobs (noise).
\end{itemize}

% =========================================================
\section{Cost Map Generation and Obstacle Inflation}
\textbf{Key function:} \texttt{generate\_cost\_map()}

The controller generates a distance transform from obstacles and then assigns traversal costs:
\begin{itemize}[leftmargin=2em]
  \item Obstacles $\rightarrow$ very high cost
  \item Unknown cells $\rightarrow$ medium cost (discourage but allow exploration if necessary)
  \item Free cells $\rightarrow$ low cost
  \item Near obstacles $\rightarrow$ inflated cost band (safety buffer)
\end{itemize}

\subsection{Adaptive Inflation}
Within \texttt{generate\_cost\_map()}, the controller samples local clearance around the robot and adjusts inflation:
\begin{itemize}[leftmargin=2em]
  \item \textbf{Tight space:} smaller inflation to fit corridors
  \item \textbf{Open space:} larger inflation to maximize safety
\end{itemize}
This adaptive behavior reduces deadlocks due to overly conservative inflation in narrow mazes.

% =========================================================
\section{Frontier-Based Exploration}
\textbf{Key functions:} \texttt{is\_frontier\_edge()}, \texttt{build\_frontier\_edge\_map()},
\texttt{find\_frontier\_centroids()}, \texttt{update\_path\_to\_frontier()},
\texttt{is\_failed\_frontier()}, \texttt{mark\_failed\_frontier()}

\subsection{Frontier Definition}
A frontier cell is a \texttt{CELL\_FREE} cell that is adjacent to at least one \texttt{CELL\_UNKNOWN} cell. These frontiers represent boundaries between known and unknown space.

\subsection{Frontier Clustering and Scoring}
\texttt{find\_frontier\_centroids()}:
\begin{enumerate}[leftmargin=2em]
  \item builds a binary frontier edge map,
  \item clusters connected frontier edge cells into components,
  \item computes centroid and size of each component,
  \item scores candidates by size and distance,
  \item sorts by score and returns the best candidates.
\end{enumerate}

\subsection{Failed Frontier Cache}
\texttt{update\_path\_to\_frontier()} tries frontiers in descending score order. If a frontier is unreachable by A*, it is stored in a TTL cache using \texttt{mark\_failed\_frontier()} to prevent repeated wasted replanning.

% =========================================================
\section{Perception (Formulas from Code)}
\textbf{Key functions:} \texttt{detect\_green\_to\_grid()}, \texttt{detect\_red\_to\_grid()}, \texttt{world\_to\_grid()}, \texttt{depth\_at\_centroid()}

This module converts camera detections (green/red pixels) into obstacle cells in the grid. The pipeline is:
\[
\text{pixel }(x,y)\ \rightarrow\ \text{bearing }\beta\ \rightarrow\ \text{world point }(w_x,w_y)\ \rightarrow\ \text{grid cell }(g_x,g_y)
\]

\subsection{Pixel $\rightarrow$ Bearing (from code)}
\textbf{Used in:} \texttt{detect\_green\_to\_grid()}, \texttt{detect\_red\_to\_grid()}
\[
\beta(x)= -\Big(x-\frac{w}{2}\Big)\cdot \frac{\mathrm{FOV}_h}{w}
\]
where $w$ is image width and $\mathrm{FOV}_h$ is horizontal FOV (e.g., 1.04 rad in code).

\subsection{Bearing $\rightarrow$ World Projection (from code)}
\textbf{Used in:} \texttt{detect\_green\_to\_grid()}, \texttt{detect\_red\_to\_grid()}
\[
\alpha = \theta + \beta
\]
\[
w_x = r_x + d\cos(\alpha),\qquad w_y = r_y + d\sin(\alpha)
\]
where $(r_x,r_y)$ is robot position, $\theta$ is robot yaw, and $d$ is depth (meters) scaled by \texttt{distance\_scale}.

\subsection{RGB Pixel $\rightarrow$ Depth Pixel Mapping (from code)}
\textbf{Used in:} \texttt{detect\_green\_to\_grid()}, \texttt{detect\_red\_to\_grid()}
\[
x_d=\mathrm{round}\Big(x\cdot \frac{d_w-1}{w-1}\Big),\qquad
y_d=\mathrm{round}\Big(y\cdot \frac{d_h-1}{h-1}\Big)
\]
Then depth is taken robustly as the \textbf{median} of a $3\times3$ neighborhood around $(x_d,y_d)$ (as implemented in code).

\subsection{World $\rightarrow$ Grid (from code)}
\textbf{Function:} \texttt{world\_to\_grid()}
\[
g_x=\Big\lfloor \frac{w_x}{\mathrm{RES}}+\frac{N}{2}\Big\rfloor,\qquad
g_y=\Big\lfloor \frac{w_y}{\mathrm{RES}}+\frac{N}{2}\Big\rfloor
\]
where $\mathrm{RES}=\texttt{GRID\_RESOLUTION}$ and $N=\texttt{GRID\_SIZE}$.

\subsection{Red Wall Segment Marking (from code)}
\textbf{Function:} \texttt{detect\_red\_to\_grid()}
After computing the red hit point $(h_x,h_y)$ at depth $d$, the code draws an \textbf{inward} obstacle segment of length $L=\texttt{segment\_len\_m}$:
\[
L'=\min(L,d)
\]
\[
(s_x,s_y)=(h_x,h_y)-L'(\cos\alpha,\sin\alpha)
\]
The segment is rasterized on the grid (Bresenham-like loop) and each traversed cell is marked as \texttt{CELL\_OBSTACLE}.

% =========================================================
\section{Pillar Detection and Target Marking}
\textbf{Key functions:} \texttt{find\_color\_centroid()}, \texttt{depth\_at\_centroid()},
\texttt{clear\_all\_pillars\_of\_type()} \\
\textbf{Main logic location:} \texttt{STATE\_MARKING\_PILLARS} (inside the state machine)

\begin{enumerate}[leftmargin=2em]
  \item detect and navigate to \textbf{blue} pillar
  \item then detect and navigate to \textbf{yellow} pillar
\end{enumerate}

\subsection{Marking Procedure}
In \texttt{STATE\_MARKING\_PILLARS}, the robot:
\begin{enumerate}[leftmargin=2em]
  \item stops motion and resets controllers,
  \item computes the color centroid in the RGB camera,
  \item computes depth at the centroid,
  \item aligns the blob to the image center (optional, skipped if too close),
  \item collects multiple world coordinate samples,
  \item averages samples and converts to grid,
  \item clears previous pillar marks and paints a disk of pillar cells.
\end{enumerate}

\subsection{Depth-Failure}
If depth is invalid (often when too close), the controller can mark a pillar once using a constant fallback distance (\texttt{FALLBACK\_DIST\_M}) and exit marking to avoid deadlock.

% =========================================================
\section{Navigation to Targets and Waypoint Following}
\textbf{Key functions:} \texttt{ensure\_path\_to\_goal()}, \texttt{follow\_path()},
\texttt{diff\_drive()}, \texttt{pid\_rotate()}, \texttt{pid\_translate()}, \texttt{apply\_slew\_limiter()}

\subsection{Path Planning to Goal}
\texttt{ensure\_path\_to\_goal()} plans or replans periodically from the robot's current cell to the goal cell. Replanning is throttled by \texttt{NAV\_REPLAN\_EVERY}.

\subsection{Waypoint Execution}
\texttt{follow\_path()} steps through \texttt{GridPath}:
\begin{itemize}[leftmargin=2em]
  \item skips waypoints that are too close,
  \item sends current waypoint into \texttt{diff\_drive()},
  \item advances to next waypoint when reached.
\end{itemize}

\subsection{Differential Drive with PID and Hysteresis}
\texttt{diff\_drive()} computes heading error and distance error to the waypoint:
\begin{itemize}[leftmargin=2em]
  \item rotates in place until aligned (hysteresis to prevent oscillation),
  \item translates forward using a velocity PID,
  \item smooths acceleration with \texttt{apply\_slew\_limiter()}.
\end{itemize}

% =========================================================
\section{Safety Manager and Recovery Behavior}
\textbf{Key functions:} \texttt{safety\_update\_per\_tick()}, \texttt{safety\_read\_ir()},
\texttt{safety\_mark\_front\_obstacle()}, \texttt{stop\_all\_motors()}, \texttt{set\_wheel\_speeds()}

\subsection{Safety Inputs}
The safety module monitors:
\begin{itemize}[leftmargin=2em]
  \item IR distances (\texttt{fl\_range}, \texttt{fr\_range})
  \item robot pitch angle (prevents wall-climb and unstable mapping)
\end{itemize}

\subsection{Safety Outputs}
Safety produces:
\begin{itemize}[leftmargin=2em]
  \item scaling factors for linear and angular velocity (\texttt{v\_scale}, \texttt{w\_scale})
  \item hard stop override (\texttt{hard\_stop}) to bypass navigation
  \item recovery behavior (backup, turn away, wait for clearance)
\end{itemize}

\subsection{Recovery State Machine}
When critical conditions occur, the robot:
\begin{enumerate}[leftmargin=2em]
  \item stops
  \item reverses for \texttt{BACKUP\_TIME\_SEC}
  \item turns away from the obstacle side
  \item waits until \texttt{FRONT\_CLEAR\_RESUME\_M} is satisfied
\end{enumerate}
Additionally, \texttt{safety\_mark\_front\_obstacle()} can paint virtual obstacles in front of the robot to prevent re-entering the same collision region.

% =========================================================
\section{Visualization and Debugging Display}
\textbf{Key functions:} \texttt{render\_display\_with\_path()}, \texttt{draw\_path\_on\_display()},
\texttt{display\_ir\_status()}, \texttt{frontier\_score\_to\_color()}, \texttt{cost\_to\_color()}

The display shows:
\begin{itemize}[leftmargin=2em]
  \item occupancy grid and pillar markers
  \item cost map heat visualization (inflation)
  \item current planned path and active waypoint
  \item detected frontiers (scored and highlighted)
  \item safety state and IR distances
\end{itemize}

This visualization is essential for interpreting navigation failures and tuning thresholds.

\end{document}